\documentclass[11pt]{article}
\usepackage[margin=1in]{geometry}
\usepackage{amsmath,amssymb,amsthm}
\title{Disproof of Conjecture 00000003836}
\author{earthking11}
\date{October 6, 2026}
\newtheorem{proposition}{Proposition}
\begin{document}
\maketitle
\section*{Pinned convention}
The classical signature algorithm encodes a reading word by signs and
repeatedly deletes adjacent $+-$ pairs.  Equivalently, scan left to right with
a stack: a minus sign cancels a plus sign at the top of the stack; otherwise
the new sign is retained.  The conjecture says that the number of remaining
plus signs is invariant under every cyclic rotation.

\section*{Two-letter counterexample}
Take
\[
 w=(+,-),\qquad \operatorname{rot}(w)=(-,+).
\]
For $w$, the unique adjacent pair is $+-$ and cancels.  Hence the reduced word
is empty and its plus count is $0$.  For the rotation, no $+-$ pair occurs, so
the reduced word is $(-,+)$ and its plus count is $1$.

\begin{proposition}
The plus count of the signature output is not invariant under cyclic rotation.
\end{proposition}
\begin{proof}
The two explicit reductions above give $0\ne1$.
\end{proof}

This refutes the first conjunct, and therefore the filed conjunction, without
needing any interpretation of the separately mentioned string parameter.
If an author uses the opposite cancellation convention ($-+$ instead of
$+-$), the very same two cyclic rotations exchange their roles and still give
counts $0$ and $1$.

The accompanying Lean file implements the stack algorithm and kernel-checks
all of these finite evaluations.
\end{document}
